\section{Experiments and Demonstrations}
\label{sec:experiments}

The experiments will combine unit-level validation, analytical parity
checks, and end-to-end demonstrations. The current test suite already
exercises most of these cases; the paper version should convert the strongest
ones into reproducible plots, tables, and ablations.

\subsection{Analytical Validation}

The first group of experiments will verify limiting cases where closed-form
or high-accuracy reference solutions are available.

\begin{table}[t]
\centering
\caption{Planned validation experiments for the conference submission.}
\label{tab:experiments}
\small
\begin{tabular}{@{}p{0.27\linewidth}p{0.35\linewidth}p{0.27\linewidth}@{}}
\toprule
Experiment & Purpose & Metric \\
\midrule
Two-body orbit & Verify chief RK4 propagation without perturbations & Relative
specific-energy drift over one orbit \\
J2 orbit & Verify secular nodal regression & RAAN drift error after one day \\
CW relative motion & Check LVLH per-body forcing in the small-offset circular
limit & Position and velocity error vs. CW solution \\
Torque-free rigid body & Check MuJoCo free-joint attitude conventions and
full-inertia handling & Energy, angular momentum, and SciPy Euler-equation error \\
Gyrostat and CMG & Validate internal momentum storage and reaction torques &
Angular-rate and momentum consistency \\
Gravity-gradient torque & Validate rotational environmental torque & Torque
direction and magnitude vs. analytical model \\
\bottomrule
\end{tabular}
\end{table}

For the CW test, a free body is initialized with meter-scale LVLH offsets and
centimeter-per-second relative velocities in a 400 km circular orbit. The
simulated relative trajectory will be compared against the classical
Clohessy-Wiltshire solution. This experiment is important because it tests
the full rotating-frame force computation while still admitting an analytical
reference.

For attitude dynamics, the paper will compare torque-free MuJoCo rollouts
against Euler-equation integration for diagonal and non-diagonal inertia
tensors. This case is also a useful place to document MuJoCo frame
conventions: free-joint angular velocity is stored in the body frame, while
\texttt{xfrc\_applied} torques are world-frame torques.

\subsection{Environmental and Actuator Case Studies}

The second group of experiments will demonstrate the simulator's spacecraft
specific features. Candidate studies include:

\begin{itemize}
  \item Drag and solar-radiation-pressure loads on configured flat plates,
  including eclipse transitions and projected-area effects.
  \item Magnetic residual-dipole and magnetorquer torques along a LEO orbit.
  \item Reaction-wheel spin-up, speed saturation, and gyrostat coupling.
  \item Thruster force/torque allocation and the resulting feedback into the
  propagated chief orbit.
  \item Single-gimbal CMG momentum and torque envelopes for a spacecraft
  attitude maneuver.
\end{itemize}

The current ISS analysis scripts provide a useful high-level demonstration
target. For the paper, these scripts can be turned into a more general
spacecraft case study with a bounded slew maneuver, environmental torque
accumulation, and actuator limits reported in SI units.

\subsection{Articulation, Contact, and Robotics Workloads}

The third group of demonstrations should emphasize what is difficult to show
in a traditional orbit propagator. The existing spacecraft-arm example
simulates a free-floating base with a two-link arm in orbit, checks that
internal arm motion remains finite and does not inject external force into
the chief orbit, and reports end-effector motion in the LVLH frame. A paper
figure could show the arm trajectory, base reaction motion, and chief-orbit
energy drift.

The contact experiment should show two free bodies colliding inside the
LVLH-frame MuJoCo model. The key result is separation of concerns: MuJoCo
contact impulses change the relative states of the bodies, but those internal
contact forces do not appear in the orbital external-wrench buffer. The paper
can report contact impulse magnitude, wrench-buffer magnitude, total LVLH
momentum before and after collision, and post-collision relative drift.

\subsection{Ablations and Reporting Plan}

Each demonstration should include an ablation that turns off one coupling
term and shows the expected change. Examples include disabling J2 in the
RAAN-drift case, disabling gravity-gradient torque in the attitude case, and
disabling orbit feedback for thruster maneuvers. The final paper should
report runtime, timestep, model size, and error tolerances so readers can
distinguish numerical accuracy from modeling assumptions.
